% ============================================================
% SECCIÓN 1
% ============================================================
\section{Introducción: Antecedentes y Parámetros Sísmicos de Diseño}
\label{sec:intro}

El proyecto corresponde al análisis y diseño estructural de un edificio residencial 
ubicado en Antofagasta, Chile, ajustado a la normativa sísmica vigente.

\subsection{Antecedentes Generales}

\begin{itemize}
    \item \textbf{Ubicación:} Antofagasta, Chile.
    \item \textbf{Uso:} Habitacional (Residencial).
    \item \textbf{Categoría de Ocupación:} Categoría II --- Estructuras de ocupación 
          normal
    \item \textbf{Zonificación Sísmica:} Zona 3 
\end{itemize}

\subsection{Parámetros Sísmicos de Diseño}

Los parámetros sísmicos de diseño, determinados según NCh433~\cite{NCh433:1996} 
y el DS N°61, se resumen en la Tabla~\ref{tab:parametros_sismicos}.

\begin{table}[H]
    \centering
    \caption{Parámetros sísmicos de diseño según NCh433 y DS N°61.}
    
    \label{tab:parametros_sismicos}
    \begin{tabular}{lc}
        \toprule
        \textbf{Parámetro}                                  & \textbf{Valor} \\
        \midrule
        Zona sísmica                                        & Zona 3 \\
        Tipo de suelo                                       & Tipo A \\
        Aceleración efectiva máxima $A_0$                   & $0{,}40\,g$ \\
        Período característico del suelo $T_0$              & $0{,}15\ \text{s}$ \\
        Factor de amplificación del suelo $S$               & $0{,}9$ \\
        Categoría de Ocupación                              & II \\
        Factor de importancia $I$                           & $1{,}0$ \\
        Sistema estructural                                 & Muros de corte de hormigón armado \\
        \bottomrule
        \vspace{0.001cm}
    \end{tabular}
    \text{Fuente: Elaboración propia}
\end{table}

\subsection{Resumen de Resultados del Prediseño Sísmico}

A modo de referencia, la Tabla~\ref{tab:resumen_prediseno} presenta los principales 
resultados del prediseño sísmico, los cuales serán contrastados con los resultados 
del modelo en secciones posteriores.

\begin{table}[H]
    \centering
    \caption{Resultados del prediseño sísmico.}
    \label{tab:resumen_prediseno}
    \begin{tabular}{lc}
        \toprule
        \textbf{Parámetro}                                  & \textbf{Valor} \\
        \midrule
        Peso sísmico del edificio $P$                       & $8{.}766{,}65\ \text{tonf}$ \\
        Período aproximado $T$                              & $1{,}05\ \text{s}$ \\
        Coeficiente sísmico calculado $C$                   & $0{,}0269$ \\
        Coeficiente sísmico utilizado $C_{\min}$            & $0{,}10$ \\
        Corte basal mínimo $Q_{\min}$                       & $526\ \text{tonf}$ \\
        Corte basal máximo $Q_{\max}$                       & $1{.}104{,}6\ \text{tonf}$ \\
        Corte basal de prediseño $Q$                        & $526\ \text{tonf}$ \\
        \bottomrule
        \vspace{0.001cm}
    \end{tabular}
    
    \text{Fuente: Elaboración propia}
\end{table}